\chapter{Die Dynamische Schwere-Trägheits-Theorie (DSTT)}

\section{Einleitung}

Die Dynamische Schwere-Trägheits-Theorie (DSTT) ist eine eigenständige Theorie der Bewegung unter dem Einfluss einer Zentralkraft. Sie entstand unabhängig von den Weber-Kräften – und doch zeigt sich, dass ihre Struktur in der WG und WED bereits enthalten ist.

Die DSTT basiert auf einer fundamentalen Erkenntnis:

\textbf{Das Äquivalenzprinzip ist falsch.}

Schwere und Trägheit sind nicht äquivalent. Sie sind nur numerisch nahe.

\section{Die Axiome der DSTT}

Die DSTT basiert auf vier Axiomen:

\subsection{Axiom 1: Trennung von Ursache und Wirkung}

Schwere (Ursache) und Trägheit (Wirkung) sind fundamental verschieden.

Die Schwere wirkt radial zwischen zwei Körpern. Es gibt keine zusätzlichen Felder, keine Raumzeitkrümmung – nur diese eine radiale Wechselwirkung.

\subsection{Axiom 2: Anisotrope Trägheit}

Die Trägheit zerfällt in zwei Komponenten: eine radiale und eine tangentiale Antwort.

Die Schwerewirkung teilt sich auf in:
\begin{itemize}
\item einen radialen Anteil \( (1-\beta) F_G \)
\item einen tangentialen Anteil \( \beta F_G \)
\end{itemize}

\subsection{Axiom 3: Dynamische Aufteilung}

Ein Zustandsparameter \( \beta(t) \in [0,1] \) bestimmt das Verhältnis.

Er ist definiert als das Verhältnis von tangentialer zu Gesamtenergie:

\[
\beta_{\text{DSTT}}(t) = \frac{E_{\text{tan}}(t)}{E(t)}
= \frac{\frac{1}{2} m r^2 \dot{\phi}^2}{\frac{1}{2} m \dot{r}^2 + \frac{1}{2} m r^2 \dot{\phi}^2}
\]

\subsection{Axiom 4: Monotonie}

Die tangentiale Komponente nimmt ständig zu:

\[
\dot{\beta}_{\text{DSTT}}(t) > 0
\]

Diese Monotonie ist eine direkte Konsequenz der Energieflüsse in der Weber-Gravitation.

\section{Die Konsequenzen der DSTT}

\subsection{Spiralbahnen statt Ellipsen}

Aus den Axiomen der DSTT folgt durch Eliminierung der Zeit:

\[
\frac{dr}{d\phi} = \frac{r}{B}, \quad B = \frac{2E}{F_G}
\]

Die Lösung ist eine logarithmische Spirale:

\[
r(\phi) = K e^{\phi/B}
\]

Kreisbahnen und Ellipsen – die Grundbausteine der newtonschen Mechanik und der ART – existieren nicht exakt. Sie sind Näherungen für kurze Zeiträume.

\subsection{Der intrinsische Zeitpfeil}

Die Monotonie \( \dot{\beta} > 0 \) bricht die Zeitumkehrinvarianz.

Die Gravitation produziert Entropie:

\[
\dot{S}_{\text{Grav}} > 0
\]

Dies liefert einen intrinsischen kosmologischen Zeitpfeil, der unabhängig von der Materie existiert.

\subsection{Die Widerlegung des Äquivalenzprinzips}

Die DSTT zeigt: Schwere und Trägheit sind nicht äquivalent.

Die Trägheit ist anisotrop – sie antwortet unterschiedlich auf radiale und tangentiale Kräfte.

Das Äquivalenzprinzip ist eine Näherung, die nur für kurze Zeitskalen und kleine Exzentrizitäten gilt.

\section{Die DSTT in der WG und WED}

Die DSTT-Struktur ist in der WG und WED bereits enthalten.

Die Weber-Kräfte haben dieselbe Aufteilung in radiale und tangentiale Komponenten.

Die DSTT ist die energetische Interpretation der Weber-Kräfte.

Die Identitätsbeziehung lautet:

\[
\frac{|F_{\text{WDBT}}^{\text{tan}}|}{|F_{\text{WDBT}}^{\text{rad}}| + |F_{\text{WDBT}}^{\text{tan}}|}
\equiv \beta_{\text{DSTT}}(t)
\]

Diese Identität gilt exakt.

\section{Zusammenfassung}

\begin{itemize}
\item Die DSTT ist eine eigenständige Theorie der Bewegung.
\item Sie zeigt: Das Äquivalenzprinzip ist falsch.
\item Sie zeigt: Spiralbahnen sind fundamental, Ellipsen sind Näherungen.
\item Sie zeigt: Die Gravitation produziert Entropie und einen Zeitpfeil.
\item Die DSTT-Struktur ist in der WG und WED enthalten.
\item Die DSTT ist die energetische Interpretation der Weber-Kräfte.
\end{itemize}